\chapter{Trend Following}\label{ch:s1:trend-following}

Take every liquid futures market, go long those that rose over the last year and short those that fell, and size each position so that it carries the
same risk. Hurst, Ooi and Pedersen ran that rule across global markets back to 1880 and found positive average returns in every decade, low correlation
with stocks and bonds, and good performance in eight of the ten worst drawdowns of a 60/40 portfolio. It is the simplest systematic strategy there is,
and one of the few with a century of evidence behind it. On this chapter's synthetic universe of forty futures the twelve-month rule earns a Sharpe ratio
of 0.66 after costs over thirty years, and in the two large stock crashes the trend book gains 21\% and 39\% while equities lose 24\% and 35\%. The build
is two components: \texttt{firm.synthfut}, the synthetic futures universe on which this and the next six chapters run, and \texttt{firm.trendfollow}.

\begin{definition}[Trend following, time-series momentum]\label{def:s1:trend-following:trend}
\emph{Trend following}\index{trend following} is trading a market in the direction of its own recent price move, long after rises and short after
falls, across many markets at once. \emph{Time-series momentum}\index{time-series momentum} is its simplest rule: the position in each market has the
sign of that market's own past excess return over a lookback of $L$ days, whatever other markets did; \omterm{def:s1:momentum:cross}{cross-sectional momentum} (chapter~5) instead ranks
markets against each other.
\end{definition}

\section{Signal families: returns, moving averages, breakouts}

Write $r_{it}$ for the daily excess return of market $i$, the return of a position rolled in the nearby futures contract, and $P_{it} =
\sum_{s \le t} r_{is}$ for its cumulative sum, a log price that ignores the roll gaps (Book~7, chapter~2's continuous futures series). Three families of
rules turn $P$ into a direction, each with a speed.
\begin{itemize}
\item \textbf{Past return.} $s_{it} = \operatorname{sign}(P_{it} - P_{i,t-L})$: \omterm{def:s1:trend-following:trend}{time-series momentum} with lookback $L$.
\item \textbf{Moving-average crossover.} $s_{it} = \operatorname{sign}(\bar P^{(n)}_{it} - \bar P^{(m)}_{it})$, the short average of $P$ over $n$ days
against the long one over $m > n$ (Book~7, chapter~7).
\item \textbf{Breakout.} The position turns long when $P$ reaches its highest level of the last $L$ days, short at its lowest, and is kept otherwise.
\end{itemize}

\begin{definition}[Breakout rule]\label{def:s1:trend-following:breakout}
A \emph{breakout rule}\index{breakout rule} takes a long position when a market's price exceeds its highest level over the last $L$ days and a short
position when it falls below its lowest, and holds the position until the opposite breakout; its channel of highs and lows is also called a Donchian
channel.
\end{definition}

All three measure the same thing, the sign of a smoothed recent slope, with different weights on the past: the past return weighs the last $L$ days
equally, a crossover weighs recent days more (a triangle of weights), and a breakout reacts only to new extremes, which makes it the slowest to trade.
\Cref{lst:s1:trend-following:signals} computes them, each from data up to the close of day $t$ for a position held over day $t + 1$.

\begin{proposition}[Trend profits are convex in the move]\label{prop:s1:trend-following:convex}
Over a window of $T$ days, a position equal to the cumulative return since the window's start, $w_t = \sum_{s < t} r_s$, earns
\[
\sum_{t=1}^{T} w_t r_t = \tfrac12\Big(\sum_{t=1}^{T} r_t\Big)^2 - \tfrac12 \sum_{t=1}^{T} r_t^2 .
\]
\end{proposition}
\begin{proof}
$\big(\sum_t r_t\big)^2 = \sum_t r_t^2 + 2\sum_{s < t} r_s r_t$, and $\sum_t w_t r_t = \sum_{s < t} r_s r_t$.
\end{proof}

The trend follower is long the square of the window's move and short its realised variance: it gains from large moves in either direction and loses in
choppy markets where the path wanders without going anywhere. That is the shape of a long straddle, and it is why the strategy's returns against the
stock market form a smile (\cref{fig:s1:trend-following:smile}). The sign rules are nonlinear versions of the same bet. Expected profits come from two
places: the square of any persistent drift, and positive autocorrelation of returns. In the synthetic universe both come from a planted drift.

\section{The synthetic futures universe}

\texttt{firm.synthfut} (\cref{lst:s1:trend-following:synthfut}) simulates thirty years of daily excess returns on forty markets, ten in each of four
classes: equity indices, government bonds, currencies and commodities, with annual volatilities of 17.3\%, 6.9\%, 9.4\% and 27.1\% as realised. Each
market has a drift that follows an AR(1) process with a half-life of a year and a standard deviation of half its volatility: the trend, planted. Each
also has a carry, of which half is earned as expected return (chapter~20 trades it). Volatility clusters within each class; shocks are Student-$t$ with a
factor common to the class. Three stock crashes, starting in years 7.5, 18 and 27 and lasting a hundred days each, add a steady drift of $-30\%$ to equities, $-10\%$ to
commodities and $+3\%$ to bonds over the episode, make high-carry currencies fall, and raise volatility by half.

The universe is built for the whole run of futures chapters: its truth (drifts, carries, volatilities, crash dates) is kept, and \texttt{curve} and
\texttt{contracts} give the log futures price for any delivery, so that chapter~21 can trade calendar spreads on its curves and chapter~23 can find its
seasonal commodities.

\section{A portfolio of trends}

Each market's position is its signal times a notional that targets 40\% annual volatility divided by the number of markets, from an exponentially
weighted volatility forecast with a centre of mass of sixty days (\cref{lst:s1:trend-following:book}); costs are 2, 1, 2 and 4 basis points per unit of
notional traded for equities, bonds, currencies and commodities (assumed). The book is measured from the second year, when every signal is defined.

\begin{center}
\small
\begin{tabular}{@{}llccc@{}}
\toprule
family & speed & Sharpe ratio after costs & before costs & turnover (times gross a year)\\
\midrule
past return & 1 month & $-0.00$ & 0.32 & 49\\
 & 3 months & 0.26 & 0.46 & 29\\
 & 12 months & 0.66 & 0.76 & 15\\
moving averages & 5/20 days & 0.10 & 0.32 & 33\\
 & 20/100 days & 0.43 & 0.48 & 9\\
 & 50/250 days & 0.72 & 0.76 & 5\\
breakout & 20 days & 0.17 & 0.29 & 18\\
 & 55 days & 0.34 & 0.39 & 8\\
 & 250 days & 0.78 & 0.80 & 3\\
\bottomrule
\end{tabular}
\end{center}

Within each family the slow rule wins, because the planted drift is slow; across families at the same speed the results are close, because the rules
measure the same slope. Costs separate the fast rules: the one-month rule turns its gross over 49 times a year and its Sharpe ratio of 0.32 before costs
becomes nothing after them. By lookback (\cref{fig:s1:trend-following:lookbacks}) the Sharpe ratio rises to six months and is flat beyond: a longer
lookback measures the drift with less noise but reacts later to its changes, and at a half-life of a year the two balance over a broad range. Moskowitz,
Ooi and Pedersen found the same range in 58 real markets: persistence for one to twelve months, partially reversing beyond.

\begin{omfigure}
\begin{tikzpicture}
\begin{axis}[width=10cm,height=5cm,xlabel={\small lookback of the past-return rule (months)},ylabel={\small Sharpe ratio after costs},
  tick label style={font=\footnotesize},xmin=0,xmax=25,ymin=-0.1,ymax=0.8,grid=major,grid style={gray!25},table/col sep=comma]
\addplot[omThm,very thick,mark=*] table[x=months,y=sr]{figdata/strategies-1/19-trend-following/lookbacks.csv};
\addplot[black!40,dashed] coordinates {(0,0) (25,0)};
\end{axis}
\end{tikzpicture}

\omcaption{\omterm{def:s1:trend-following:trend}{Time-series momentum} on the synthetic universe of forty futures, thirty years: the Sharpe ratio after costs of the volatility-scaled
sign-of-past-return rule by lookback. The planted drifts have a half-life of a year. Data: \texttt{s1\_trend.lookbacks}.}\label{fig:s1:trend-following:lookbacks}
\end{omfigure}

A practitioner does not know the speed of the trends to come, so trend books blend speeds. The blend of the three past-return rules, averaging their
positions, earns 0.40; the blend of all nine rules, 0.55. Both are below the best single rule, which is what hindsight costs: the best rule is known only
after the fact, and choosing it among nine is the selection that Book~7, chapter~20's overfitting measures penalise. Blending is insurance against the
speed of the next trend being different from the last.

\section{Volatility scaling}

\begin{definition}[Volatility targeting]\label{def:s1:trend-following:voltarget}
\emph{Volatility targeting}\index{volatility targeting} sizes a position, or a whole portfolio, so that its forecast volatility equals a fixed target:
the notional is the target divided by the forecast, recomputed as the forecast changes, usually with a cap on leverage.
\end{definition}

Trend books use it at two levels. Per market, it makes a bond future and a natural-gas future carry the same risk, so that the book is diversified in
risk and not dominated by the most volatile markets. Replacing the three past-return rules' volatility-scaled positions with equal notional positions
drops the blend's Sharpe ratio from 0.40 to 0.31: the commodities dominate the equal-notional book and the bonds hardly count. Baltas and Kosowski
show that the choice of volatility estimator matters too, mainly through turnover.

At the portfolio level, targeting 10\% raises the blend's realised volatility from 7.7\% to 10.2\% and lowers its Sharpe ratio from 0.40 to 0.34; the
largest drawdown, 3.7 annual standard deviations untargeted, is 3.8 targeted. In this model the book's volatility is mostly a measure of how many markets
trend together: when they disagree, the book's volatility falls and the target levers it up, just when its edge is small. Portfolio targeting is a risk
decision, not an alpha, and chapter~26 measures when it helps.

\section{The long record and crisis alpha}

\begin{definition}[Crisis alpha]\label{def:s1:trend-following:crisis}
\emph{Crisis alpha}\index{crisis alpha} is a strategy's tendency to earn positive returns during prolonged declines of the stock market, measured as its
return over the market's largest drawdowns; for \omterm{def:s1:trend-following:trend}{trend following} it comes from being short equities and positioned for the flight to quality once the
decline has lasted longer than the signal's lookback.
\end{definition}

The synthetic crashes show the mechanism (\cref{fig:s1:trend-following:paths}). The first, in year 7.5, happened to be small (equities lost 6.8\% over
its hundred days, the planted fall offset by the noise) and the trend book, at 10\% volatility, gained 2.3\%. In the second, equities fell 24.3\% and the
book gained 21.1\%; in the third, 35.1\% and 39.0\%. The gains come once the declines are under way: the signals turn short equities and long bonds a few
weeks in, and the steady fall does the rest. A crash that happens in a day, like 19 October 1987 (Book~6, chapter~22), gives the signals no time and
earns nothing.

\begin{omfigure}
\begin{tikzpicture}
\begin{axis}[width=11cm,height=5.4cm,xlabel={\small year},ylabel={\small cumulative return (\%)},tick label style={font=\footnotesize},
  xmin=1,xmax=30,ymin=-20,ymax=160,grid=major,grid style={gray!25},legend columns=2,legend style={font=\scriptsize,at={(0.5,-0.3)},
  anchor=north,draw=none,fill=none,/tikz/every even column/.append style={column sep=8pt}},table/col sep=comma]
\fill[omExo!15] (axis cs:7.5,-20) rectangle (axis cs:7.897,160);
\fill[omExo!15] (axis cs:18,-20) rectangle (axis cs:18.397,160);
\fill[omExo!15] (axis cs:27,-20) rectangle (axis cs:27.397,160);
\addplot[omThm,thick] table[x=year,y=trend]{figdata/strategies-1/19-trend-following/paths.csv};
\addplot[omDef,thick] table[x=year,y=equity]{figdata/strategies-1/19-trend-following/paths.csv};
\legend{trend book (10\% volatility target),equity class (mean of ten indices)}
\end{axis}
\end{tikzpicture}

\omcaption{The blended trend book, targeted at 10\% volatility, and the mean of the ten equity index futures on the synthetic universe, as cumulative sums
of daily excess returns; the shaded bands are the three planted stock crashes. Data: \texttt{s1\_trend.paths}.}\label{fig:s1:trend-following:paths}
\end{omfigure}

Quarter by quarter, the book's returns against the equity class's form the smile of the convexity proposition
(\cref{fig:s1:trend-following:smile}): a quadratic fitted to the 115 quarters has a curvature of 2.63, and the worst equity quarter ($-36.2\%$) was the
book's best ($+31.6\%$). Across all quarters the correlation is $-0.09$: the trend book is not a short on stocks, but a long position in large moves.

\begin{omfigure}
\begin{tikzpicture}
\begin{axis}[width=10cm,height=5.6cm,xlabel={\small equity class return in the quarter (\%)},ylabel={\small trend book return (\%)},
  tick label style={font=\footnotesize},xmin=-40,xmax=30,ymin=-20,ymax=40,grid=major,grid style={gray!25},legend pos=north east,
  legend style={font=\scriptsize},table/col sep=comma]
\addplot[only marks,mark=*,mark size=1.2pt,omDef] table[x=equity,y=trend]{figdata/strategies-1/19-trend-following/smile.csv};
\addplot[omThm,very thick] table[x=equity,y=trend]{figdata/strategies-1/19-trend-following/smile_fit.csv};
\legend{quarters,quadratic fit}
\end{axis}
\end{tikzpicture}

\omcaption{The trend smile: the blended trend book's quarterly return (10\% volatility target) against the equity class's, 115 quarters of the synthetic
universe, with a fitted quadratic (curvature 2.63). Data: \texttt{s1\_trend.smile}.}\label{fig:s1:trend-following:smile}
\end{omfigure}

The public record is longer than any synthetic one. Moskowitz, Ooi and Pedersen found a diversified \omterm{def:s1:trend-following:trend}{time-series momentum} portfolio with substantial
abnormal returns, little exposure to standard factors and its best performance in extreme markets; Hurst, Ooi and Pedersen extended the evidence to 1880,
with positive average returns in every decade and good performance in eight of the ten largest crisis periods.

One market can be tested on real data here: NYMEX WTI crude oil, from the EIA's daily settlements of the first two contracts (Book~3, chapter~2),
rolled five days before expiry into a continuous series. The twelve-month rule, at a 40\% volatility target, earned 9.0\% a year from 1986 to 2024 at a
Sharpe ratio of 0.22, long on 54\% of days; holding the future long earned a Sharpe ratio of 0.04. One market is a small sample (a Sharpe ratio measured
over 38 years has a standard error of about 0.16), and \omterm{def:s1:trend-following:trend}{trend following}'s case rests on the many markets that diversify it.

\section{The recent decade}

A long record does not guarantee the next decade. Baltas and Kosowski study the underperformance of \omterm{def:s1:trend-following:trend}{time-series momentum} after the global financial
crisis and relate it to the level of pairwise correlations between markets: when every market moves with every other, the book holds one position forty
times. The synthetic universe gives the other half of the answer, sampling noise. The blended book, whose planted edge is the same throughout, had
Sharpe ratios of 0.56, 0.24 and 0.34 in its three decades. A decade's Sharpe ratio has a standard error of about $1/\sqrt{10} = 0.32$, so a strategy
with a true Sharpe ratio of 0.4 will show anything from 0.08 to 0.72 in most decades. A bad decade is weak evidence that the edge has gone; a
mechanism that explains why it would have gone (correlations, crowding, costs) is stronger.

\section{Strategy files}

\begin{strategyfile}{Time-series momentum}\label{strat:s1:trend-following:tsmom}
\sfield{Who pays you, and why} Hedgers who trade against trends, and investors who underreact to news and overreact later; speculators profit at hedgers'
expense in Moskowitz, Ooi and Pedersen's data.
\sfield{Instruments and venues} Liquid futures in equity indices, bonds, currencies and commodities.
\sfield{Signal} The sign of each market's past 12-month excess return.
\sfield{Sizing and execution} Volatility-scaled positions, rebalanced daily or weekly; futures rolled before expiry.
\sfield{Costs} Low: a turnover of about fifteen times gross a year at a twelve-month lookback.
\sfield{How it dies} Correlated markets; choppy markets without persistent moves; crowding.
\sfield{Horizon, capacity, infrastructure} Months; large capacity in the liquid futures; a roll and margin operation.
\sfield{Backtest honestly} Rolled contracts with their costs; volatility forecasts known at the time; lookback not chosen in hindsight.
\sfield{Sources} Moskowitz, Ooi and Pedersen (2012); this chapter: 0.66 after costs on the synthetic universe.
\end{strategyfile}

\begin{strategyfile}{Moving-average crossover portfolio}\label{strat:s1:trend-following:ma}
\sfield{Who pays you, and why} As for \omterm{def:s1:trend-following:trend}{time-series momentum}.
\sfield{Instruments and venues} As for \omterm{def:s1:trend-following:trend}{time-series momentum}.
\sfield{Signal} The sign of a short minus a long moving average of each market's price, for example 50/250 days.
\sfield{Sizing and execution} Volatility-scaled; the crossover trades less often than the past-return rule at a similar speed.
\sfield{Costs} Lowest of the three families at the slow speeds (five times gross a year at 50/250 here).
\sfield{How it dies} As for \omterm{def:s1:trend-following:trend}{time-series momentum}.
\sfield{Horizon, capacity, infrastructure} Months; as for \omterm{def:s1:trend-following:trend}{time-series momentum}.
\sfield{Backtest honestly} The pair of windows fixed before the test, or chosen with Book~7's overfitting controls.
\sfield{Sources} Baltas and Kosowski (2020) on trading rules and turnover; this chapter: 0.72.
\end{strategyfile}

\begin{strategyfile}{Breakout system}\label{strat:s1:trend-following:breakout}
\sfield{Who pays you, and why} As for \omterm{def:s1:trend-following:trend}{time-series momentum}; breakouts also meet stop orders placed at recent highs and lows.
\sfield{Instruments and venues} Liquid futures.
\sfield{Signal} New $L$-day highs and lows (20, 55 or 250 days).
\sfield{Sizing and execution} Volatility-scaled; orders at the channel, often as stops.
\sfield{Costs} Slippage at breakouts, when many stops trigger together.
\sfield{How it dies} False breakouts in ranges; crowded stops.
\sfield{Horizon, capacity, infrastructure} Weeks to months.
\sfield{Backtest honestly} Fills at the breakout level plus slippage, not at the day's close.
\sfield{Sources} No performance figure verified; this chapter: 0.78 at 250 days, 0.17 at 20.
\end{strategyfile}

\begin{strategyfile}{Multi-speed trend blend}\label{strat:s1:trend-following:blend}
\sfield{Who pays you, and why} As for \omterm{def:s1:trend-following:trend}{time-series momentum}, at every horizon.
\sfield{Instruments and venues} Liquid futures.
\sfield{Signal} The average of rules at several speeds and families.
\sfield{Sizing and execution} Positions averaged before trading, so that offsetting rules net.
\sfield{Costs} Those of the fast components, reduced by netting.
\sfield{How it dies} As for \omterm{def:s1:trend-following:trend}{time-series momentum}.
\sfield{Horizon, capacity, infrastructure} Weeks to a year.
\sfield{Backtest honestly} The set of speeds fixed in advance; compare with the best single rule only as hindsight.
\sfield{Sources} Hurst, Ooi and Pedersen (2017); this chapter: 0.55 for nine rules, 0.40 for three speeds of one.
\end{strategyfile}

\begin{strategyfile}{Volatility-scaled trend portfolio}\label{strat:s1:trend-following:volscaled}
\sfield{Who pays you, and why} As for \omterm{def:s1:trend-following:trend}{time-series momentum}.
\sfield{Instruments and venues} Liquid futures across the four classes.
\sfield{Signal} Any trend rule.
\sfield{Sizing and execution} Each market at the same forecast volatility; the portfolio optionally at a volatility target with a leverage cap.
\sfield{Costs} Extra turnover from volatility changes; efficient estimators reduce it.
\sfield{How it dies} Volatility forecasts that lag regime changes.
\sfield{Horizon, capacity, infrastructure} As for the rule; a daily volatility pipeline.
\sfield{Backtest honestly} Forecasts from past data only; leverage caps as they would have applied.
\sfield{Sources} Moskowitz, Ooi and Pedersen (2012); Baltas and Kosowski (2020); this chapter: 0.40 scaled against 0.31 equal notional.
\end{strategyfile}

\begin{strategyfile}{Trend on WTI crude futures}\label{strat:s1:trend-following:wti}
\sfield{Who pays you, and why} Producers and consumers hedging, and the slow adjustment of oil prices to supply shocks.
\sfield{Instruments and venues} NYMEX WTI crude futures, rolled before expiry.
\sfield{Signal} The sign of the past 12-month return of the rolled series.
\sfield{Sizing and execution} A 40\% volatility target, rolled five days before expiry.
\sfield{Costs} The roll's spread and the curve's shape (contango costs a long position).
\sfield{How it dies} As one market, it is not diversified: it dies with oil's regime.
\sfield{Horizon, capacity, infrastructure} Months.
\sfield{Backtest honestly} The rolled series, not the nearby price; April 2020's negative settlement handled by rolling early.
\sfield{Sources} EIA settlements; this chapter: a Sharpe ratio of 0.22, 1986--2024.
\end{strategyfile}

\section{Tutorial: crisis alpha}\label{tut:s1:trend-following:tut}

\begin{tutorial}
\textbf{Goal.} Build the synthetic futures universe, run the trend rules at three speeds in three families, blend and scale them, and measure the book in
the planted crashes and on real crude oil. \textbf{End state:} the table and the three figures.
\begin{tutsteps}
\item \textbf{The universe}: drifts, carries, volatility states, class factors and crashes, day by day.
\omcode{code/firm/synthfut/firm_synthfut.py}{89}{103}{One day of the synthetic futures universe.}\label{lst:s1:trend-following:synthfut}
\item \textbf{The signals}: volatility forecast, past return, moving-average crossover and breakout.
\omcode{code/firm/trendfollow/firm_trendfollow.py}{24}{60}{Trend signals and the volatility forecast.}\label{lst:s1:trend-following:signals}
\item \textbf{The book}: positions, P\&L net of costs, and the portfolio volatility target.
\omcode{code/firm/trendfollow/firm_trendfollow.py}{63}{84}{Positions, P\&L and volatility targeting.}\label{lst:s1:trend-following:book}
\item \textbf{Run} \texttt{table()}, \texttt{lookbacks()}, \texttt{summary()}, \texttt{crises()}, \texttt{smile()}, \texttt{decades()} and
\texttt{wti()}, and \texttt{fig\_trend.py}.
\end{tutsteps}
\textbf{What to change next.} Shorten the drifts' half-life and find the best lookback; make a crash happen in five days instead of a hundred; replace
the sign with a continuous signal (the past return divided by its volatility) and compare turnover.
\end{tutorial}

\section{Build: futures universe and trend library}\label{bld:s1:trend-following:trend}

\begin{build}
\bfield{Purpose} \texttt{firm.synthfut}: a deterministic futures universe with planted drifts, carry, volatility clustering, crashes and seasonal curves,
for chapters 19 to 25. \texttt{firm.trendfollow}: trend signals, volatility-scaled positions, P\&L with costs and portfolio targeting.
\bfield{Interface} \texttt{FutConfig(\dots)}, \texttt{simulate\_futures(cfg)}, \texttt{season(amp, phase, t)}, \texttt{curve(F, t, taus)},
\texttt{contracts(F, t, n, every)}; \texttt{ewma\_vol}, \texttt{tsmom}, \texttt{ma\_cross}, \texttt{breakout}, \texttt{positions}, \texttt{run},
\texttt{vol\_target}.
\bfield{Rules} Signals use data up to the close of $t$ for positions held over $t + 1$; the universe keeps its truth; futures returns carry no seasonality
(the curve anticipates it).
\bfield{Acceptance tests} \texttt{code/firm/synthfut/tests/} and \texttt{code/firm/trendfollow/tests/}: shapes and determinism, class volatilities,
a crash's fall, the curve's slope equal to the carry; signals on a hand-made path, P\&L and costs by hand, the volatility forecast and target.
\bfield{Stretch} Continuous signals; an EWMA crossover family; roll schedules and contract-level P\&L.
\end{build}

\begin{omsources}
\item T.~J.~Moskowitz, Y.~H.~Ooi and L.~H.~Pedersen, ``Time series momentum'', \emph{Journal of Financial Economics} 104(2), 2012.
\item B.~Hurst, Y.~H.~Ooi and L.~H.~Pedersen, ``A century of evidence on trend-following investing'', \emph{Journal of Portfolio Management} 44(1), 2017.
\item N.~Baltas and R.~Kosowski, ``Demystifying time-series momentum strategies: volatility estimators, trading rules and pairwise correlations'', in
\emph{Market Momentum}, Wiley, 2020.
\item US Energy Information Administration, NYMEX WTI crude oil futures daily settlements, contracts 1 and 2.
\end{omsources}

\section{Exercises}

\begin{exercise}[$\star$]\label{exo:s1:trend-following:1}
An exponentially weighted variance with a centre of mass of 60 days has a decay factor $\lambda = 60/61$. What is it, and what is the weights'
half-life in days?
\end{exercise}

\begin{exercise}[$\star$]\label{exo:s1:trend-following:2}
A book of 40 markets targets 40\% volatility per market divided by 40. What notional, as a share of capital, does it hold in a commodity with 25\%
volatility, and in a bond future with 6\%?
\end{exercise}

\begin{exercise}[$\star$]\label{exo:s1:trend-following:3}
Use the convexity proposition to explain why a trend follower loses in a market that rises 5\% and falls 5\% repeatedly.
\end{exercise}

\begin{exercise}[$\star\star$]\label{exo:s1:trend-following:4}
The smile's fitted quadratic is $2.63x^2 + 0.03x - 0.01$ in quarterly returns. What does it predict for the trend book in a quarter when equities fall
30\%?
\end{exercise}

\begin{exercise}[$\star\star$]\label{exo:s1:trend-following:5}
Why does the one-month rule, with a Sharpe ratio of 0.32 before costs, earn nothing after them, while the 250-day breakout loses almost nothing to costs?
\end{exercise}

\begin{exercise}[$\star\star$]\label{exo:s1:trend-following:6}
A trend book's true Sharpe ratio is 0.4. Roughly what range of decade Sharpe ratios should you expect, and what do you conclude from a decade at 0.1?
\end{exercise}

\begin{exercise}[$\star\star\star$]\label{exo:s1:trend-following:7}
\emph{Coding.} Simulate the universe with \texttt{FutConfig(trend\_half\_life=63.0)} and rerun \texttt{lookbacks((3, 6, 12))}. Which lookback is best,
and why is it not three months?
\end{exercise}

\begin{exercise}[$\star\star\star$]\label{exo:s1:trend-following:8}
\emph{Find the flaw.} ``We tested 200 combinations of lookbacks and moving-average windows on our 40 futures and the best earned a Sharpe ratio of 1.3
after costs; we will run it.''
\end{exercise}

\section{Problem: Crisis Alpha}

\begin{problem}[{Weekend problem --- a trend book through three crashes}]\label{pb:s1:trend-following:1}
The chapter's synthetic futures universe, the WTI series and the public record.

\medskip\noindent\textbf{Part I --- Rules.}
\begin{enumerate}
\item Define \omterm{def:s1:trend-following:trend}{trend following}, \omterm{def:s1:trend-following:trend}{time-series momentum} and a \omterm{def:s1:trend-following:breakout}{breakout rule}.
\item Write the three signal families and say how they weight the past.
\item State and prove the convexity proposition.
\item What two sources of expected profit does a trend rule have?
\end{enumerate}

\medskip\noindent\textbf{Part II --- The universe and the book.}
\begin{enumerate}[resume]
\item Describe \texttt{firm.synthfut}'s drifts, carries, volatility and crashes.
\item How are positions sized, and what do costs assume?
\item Give the Sharpe ratios after costs of the nine rules.
\item Why do slow rules win here, and what does blending cost and buy?
\end{enumerate}

\medskip\noindent\textbf{Part III --- Scaling and crises.}
\begin{enumerate}[resume]
\item Define \omterm{def:s1:trend-following:voltarget}{volatility targeting}; compare volatility-scaled and equal-notional books.
\item What does portfolio targeting do to the blend, and why?
\item Define \omterm{def:s1:trend-following:crisis}{crisis alpha} and give the book's returns in the three crashes.
\item Describe the smile and its curvature.
\end{enumerate}

\medskip\noindent\textbf{Part IV --- The verdict.}
\begin{enumerate}[resume]
\item State the \emph{named result}: the trend portfolio's return in the simulated stock crashes and its Sharpe ratio across speeds.
\item What did the twelve-month rule earn on WTI, and how much should one market convince you?
\item What did Moskowitz, Ooi and Pedersen, and Hurst, Ooi and Pedersen find?
\item What do Baltas and Kosowski relate the recent underperformance to?
\item How noisy is a decade's Sharpe ratio?
\item Which crash would \omterm{def:s1:trend-following:trend}{trend following} not protect against?
\item How does this chapter's momentum differ from chapter~5's?
\item In one sentence: what does a trend follower own?
\end{enumerate}
\end{problem}

\section{Interview questions}

\begin{interviewq}[$\star$ \iqroles{researcher}]\label{iq:s1:trend-following:1}
What is the difference between time-series and \omterm{def:s1:momentum:cross}{cross-sectional momentum}?
\end{interviewq}

\begin{interviewq}[$\star\star$ \iqroles{researcher}]\label{iq:s1:trend-following:2}
Why do trend followers scale positions by volatility?
\end{interviewq}

\begin{interviewq}[$\star\star$ \iqroles{researcher, trader}]\label{iq:s1:trend-following:3}
Your trend book has lost money for two years. How do you decide whether to keep running it?
\end{interviewq}

\begin{interviewq}[$\star\star$ \iqroles{developer}]\label{iq:s1:trend-following:4}
What does a futures backtester need that an equity backtester does not?
\end{interviewq}

\begin{interviewq}[$\star\star$ \iqroles{risk}]\label{iq:s1:trend-following:5}
Is \omterm{def:s1:trend-following:trend}{trend following} a hedge for an equity portfolio? When is it not?
\end{interviewq}

\begin{interviewq}[$\star\star\star$ \iqroles{researcher}]\label{iq:s1:trend-following:6}
Returns have a constant drift $\mu$ and independent noise of variance $\sigma^2$ per day. A linear rule holds $w_t = \sum_{k=1}^{L} r_{t-k}$. Derive its
expected daily P\&L and the Sharpe ratio of that P\&L for small $\mu/\sigma$.
\end{interviewq}
